% w12_s02_ficha_recepcion_evidencia_bloqueo.tex
% IMT-221 — Semana 12, Sesión 2 (21/10/2026)
% Artefacto 1: Ficha de recepción de evidencia y bloqueo técnico
\PassOptionsToPackage{table,xcdraw}{xcolor}
\documentclass[10pt, a4paper]{article}

\usepackage[T1]{fontenc}
\usepackage[utf8]{inputenc}
\usepackage{tgpagella}
\usepackage{tgcursor}
\usepackage[spanish, es-noshorthands]{babel}
\usepackage{amsmath, amssymb}
\usepackage{booktabs}
\usepackage{tabularx}
\usepackage{xcolor}
\usepackage{tikz}
\usetikzlibrary{shapes.geometric, arrows.meta, positioning}
\usepackage[most]{tcolorbox}
\usepackage[margin=1.5cm]{geometry}
\usepackage{fancyhdr}

\setlength{\headheight}{16pt}
\setlength{\emergencystretch}{2em}
\definecolor{ucbblue}{RGB}{0, 51, 102}
\definecolor{neutralbg}{RGB}{247, 248, 250}
\definecolor{warnbg}{RGB}{254, 246, 246}

\pagestyle{fancy}
\fancyhf{}
\lhead{\textbf{IMT-221}}
\chead{Ficha de Recepción de Evidencia y Bloqueo Técnico (Gate A)}
\rhead{Semana 12 - Sesión 2}
\lfoot{Unidad 2: Amplificación de Precisión — Suplemento de Trazabilidad}
\rfoot{Página \thepage}
\renewcommand{\headrulewidth}{0.4pt}

\begin{document}

\begin{center}
    {\Large\textbf{\textcolor{ucbblue}{Ficha de Recepción de Evidencia y Dictamen de Bloqueo Técnico (Hito 2)}}} \\[0.1cm]
    {\footnotesize\textbf{Suplemento de Trazabilidad Documental (15 min) · Reutiliza criterios técnicos de \texttt{w12\_s01\_verificacion\_gate\_A.pdf}}}
\end{center}

\vspace{-0.2cm}
\begin{center}
    \begin{tikzpicture}[>=Stealth, thick, scale=0.84, transform shape]
        \node[draw=ucbblue, rectangle, rounded corners=2pt, fill=ucbblue!10, minimum height=0.8cm, text width=2.0cm, align=center, font=\scriptsize\bfseries] (n1) at (0,0) {1. Hito 1\\Trazable};
        \node[draw=ucbblue, rectangle, rounded corners=2pt, fill=ucbblue!10, minimum height=0.8cm, text width=2.1cm, align=center, font=\scriptsize\bfseries] (n2) at (2.7,0) {2. Requisito\\y Métrica};
        \node[draw=ucbblue, rectangle, rounded corners=2pt, fill=ucbblue!10, minimum height=0.8cm, text width=2.1cm, align=center, font=\scriptsize\bfseries] (n3) at (5.4,0) {3. Ganancia\\Objetivo $A_d^\star$};
        \node[draw=ucbblue, rectangle, rounded corners=2pt, fill=ucbblue!10, minimum height=0.8cm, text width=2.1cm, align=center, font=\scriptsize\bfseries] (n4) at (8.1,0) {4. Síntesis $R_G$\\y Matching};
        \node[draw=ucbblue, rectangle, rounded corners=2pt, fill=ucbblue!10, minimum height=0.8cm, text width=2.1cm, align=center, font=\scriptsize\bfseries] (n5) at (10.8,0) {5. Simulación\\LTSpice};
        \node[draw=ucbblue, rectangle, rounded corners=2pt, fill=ucbblue!10, minimum height=0.8cm, text width=2.1cm, align=center, font=\scriptsize\bfseries] (n6) at (13.5,0) {6. Fuente, IC\\y Seguridad};
        \node[draw=green!60!black, rectangle, rounded corners=2pt, fill=green!15, minimum height=0.8cm, text width=2.0cm, align=center, font=\scriptsize\bfseries] (n7) at (16.2,0) {7. Dictamen\\Gate A};
        \draw[->] (n1) -- (n2); \draw[->] (n2) -- (n3); \draw[->] (n3) -- (n4);
        \draw[->] (n4) -- (n5); \draw[->] (n5) -- (n6); \draw[->] (n6) -- (n7);
    \end{tikzpicture}
\end{center}

\vspace{-0.2cm}
\begin{tcolorbox}[colback=neutralbg, colframe=ucbblue, boxsep=2pt, top=3pt, bottom=3pt, title=\textbf{Regla de Trazabilidad e Inspección Rápida}]
    \footnotesize\raggedright
    Una colección de archivos o capturas \textbf{no equivale} a un diseño si no permite seguir el vínculo entre datos, cálculos y validación. El docente debe poder abrir la fuente exacta de cada dato. Si falta una medición base, escriba explícitamente \texttt{NO MEDIDO / EVIDENCIA NO DISPONIBLE} (prohibido extrapolar datos ajenos). \\
    \textbf{Ejemplo ilustrativo (NO son datos del curso):} \texttt{Hito 1: $V_{d,\max}=40\,\text{mV}$ en condición X, archivo medidas.csv fila 8 $\rightarrow$ $V_{o,\text{target}}=0.8\,\text{V}$ $\rightarrow$ $A_d^\star=20$ $\rightarrow$ esquema h2.asc $\rightarrow$ LTSpice común: PENDIENTE}.
\end{tcolorbox}

\section*{1. Ficha Única de Evidencia Trazable por Grupo}
\begin{table}[htbp]
    \centering
    \scriptsize
    \renewcommand{\arraystretch}{1.28}
    \begin{tabularx}{\textwidth}{|>{\raggedright\arraybackslash}p{4.6cm}|>{\raggedright\arraybackslash}p{4.2cm}|c|>{\raggedright\arraybackslash}X|}
    \hline
    \rowcolor[HTML]{EFEFEF} 
    \textbf{Campo de Evidencia Requerido} & \textbf{Valor o Referencia Declarada por el Grupo} & \textbf{Estado (V / PNV / A)} & \textbf{Archivo, Hoja, Medición o Captura y Ubicación Exacta} \\ \hline
    1. Grupo, integrantes y fecha & & $\square$ V $\square$ PNV $\square$ A & \\ \hline
    2. Hito 1: configuración efectiva (Wheatstone/NTC, Twin-T, nodos) & & $\square$ V $\square$ PNV $\square$ A & \\ \hline
    3. Hito 1: baseline real y condición de medición & & $\square$ V $\square$ PNV $\square$ A & \\ \hline
    4. Limitación observable y evidencia cuantitativa & & $\square$ V $\square$ PNV $\square$ A & \\ \hline
    5. Métrica propuesta de mejora sistémica & & $\square$ V $\square$ PNV $\square$ A & \\ \hline
    6. Diagrama propuesto Hito 2 y ubicación del In-Amp & & $\square$ V $\square$ PNV $\square$ A & \\ \hline
    7. Rango de $V_d$, $V_{cm}$, salida objetivo y frecuencia & & $\square$ V $\square$ PNV $\square$ A & \\ \hline
    8. $A_d^\star$, $R, R_G, k$, resistores disponibles y $A_{d,\mathrm{nom}}$ recalculada & & $\square$ V $\square$ PNV $\square$ A & \\ \hline
    9. Matching crítico ($R_2/R_1, R_4/R_3$) y estrategia CMRR & & $\square$ V $\square$ PNV $\square$ A & \\ \hline
    10. LTSpice: modo diferencial, modo común y headroom & & $\square$ V $\square$ PNV $\square$ A & \\ \hline
    11. IC preciso, fabricante, datasheet y pinout verificado & & $\square$ V $\square$ PNV $\square$ A & \\ \hline
    12. Fuente, nodo de referencia ($V_{REF}$) y márgenes eléctricos & & $\square$ V $\square$ PNV $\square$ A & \\ \hline
    13. BOM completa y disponibilidad física real & & $\square$ V $\square$ PNV $\square$ A & \\ \hline
    14. Siguiente medición o comprobación planificada & & $\square$ V $\square$ PNV $\square$ A & \\ \hline
    \end{tabularx}
\end{table}
\noindent{\scriptsize\textbf{Códigos de Estado:} \textbf{V} = Verificable (inspeccionado y consistente); \textbf{PNV} = Presente No Verificado (declarado sin respaldo comprobable inmediato); \textbf{A} = Ausente (\texttt{NO MEDIDO / EVIDENCIA NO DISPONIBLE}).}

\section*{2. Semáforo de Decisión Gate A y Clasificación de Faltantes}
\begin{center}
    \begin{tikzpicture}[>=Stealth, thick, scale=0.85, transform shape]
        \node[draw, rectangle, rounded corners=2pt, fill=white, text width=3.1cm, align=center, minimum height=0.9cm, font=\footnotesize\bfseries] (gate) at (0,0) {Evaluación Gate A\\(Ficha + Defensa Oral)};
        
        \node[draw=green!60!black, rectangle, fill=green!12, text width=4.3cm, align=center, minimum height=0.8cm, font=\scriptsize] (r1) at (5.2, 1.15) {\textbf{READY}\\Topología, cálculos, SPICE, IC, rieles y conexiones válidas};
        \node[draw=orange!80!black, rectangle, fill=yellow!20, text width=4.3cm, align=center, minimum height=0.8cm, font=\scriptsize] (r2) at (5.2, 0) {\textbf{MINOR CORRECTIONS}\\Diseño defendible con detalle no crítico o ajuste localizado};
        \node[draw=red!80!black, rectangle, fill=red!12, text width=4.3cm, align=center, minimum height=0.8cm, font=\scriptsize] (r3) at (5.2, -1.15) {\textbf{NOT READY}\\Falta fundamento: topología, $A_d$, SPICE, pinout, fuente o baseline};

        \node[draw=green!60!black, rectangle, fill=green!5, text width=5.2cm, align=center, minimum height=0.8cm, font=\scriptsize] (a1) at (11.4, 1.15) {Autoriza montaje aislado $\rightarrow$ \textbf{Energizar solo tras chequeo físico previo y supervisión}};
        \node[draw=orange!80!black, rectangle, fill=yellow!8, text width=5.2cm, align=center, minimum height=0.8cm, font=\scriptsize] (a2) at (11.4, 0) {Corregir ítem puntual $\rightarrow$ \textbf{Revalidar antes de autorizar energización}};
        \node[draw=red!80!black, rectangle, fill=red!5, text width=5.2cm, align=center, minimum height=0.8cm, font=\scriptsize] (a3) at (11.4, -1.15) {Permanecer en diseño/simulación $\rightarrow$ \textbf{PROHIBIDO ENERGIZAR}};

        \draw[->] (gate.east) -- ++(0.4,0) |- (r1.west);
        \draw[->] (gate.east) -- (r2.west);
        \draw[->] (gate.east) -- ++(0.4,0) |- (r3.west);
        \draw[->] (r1.east) -- (a1.west);
        \draw[->] (r2.east) -- (a2.west);
        \draw[->] (r3.east) -- (a3.west);
    \end{tikzpicture}
\end{center}

\noindent{\footnotesize\textbf{Distinción Obligatoria de Faltantes:}
\begin{itemize}
    \setlength{\itemsep}{1pt}
    \item \textbf{Bloqueo técnico crítico ($\rightarrow$ \texttt{NOT READY}):} Ganancia equivocada o sin justificar desde Hito 1; realimentación incorrecta; simulación contradictoria sin diagnóstico; voltajes fuera de rango en datasheet; tierra de instrumento peligrosa frente a $V_{REF}$; ausencia de evidencia mínima para fijar el objetivo. \textit{Un archivo faltante que impide verificar un ensayo técnico es un bloqueo crítico, no una brecha menor.}
    \item \textbf{Brecha documental menor ($\rightarrow$ \texttt{MINOR CORRECTIONS}):} Página mal numerada, captura pendiente de pegar en el documento o formato incompleto, \textbf{únicamente si el docente ya inspeccionó y validó en vivo la evidencia técnica subyacente}.
\end{itemize}}

\begin{tcolorbox}[colback=warnbg, colframe=red!75!black, boxsep=2pt, top=3pt, bottom=3pt, title=\textbf{Acta Mínima de Dictamen Gate A y Autorización de Energización}]
    \footnotesize
    \textbf{Grupo:} \rule{2.8cm}{0.4pt} \hfill \textbf{Dictamen:} $\square$ \textbf{READY} \quad $\square$ \textbf{MINOR CORRECTIONS} \quad $\square$ \textbf{NOT READY} \\[0.15cm]
    \textbf{Evidencia técnica efectivamente inspeccionada:} \hrulefill \\[0.15cm]
    \textbf{Bloqueo prioritario causal (Uno solo, raíz del problema):} \hrulefill \\[0.15cm]
    \textbf{Próxima acción específica y verificable:} \hrulefill \\[0.15cm]
    \textbf{Responsable:} \rule{4.5cm}{0.4pt} \hfill \textbf{Fecha/Hora de reverificación:} \rule{4.0cm}{0.4pt} \\[0.15cm]
    \textbf{Energización Física Autorizada (Por defecto = NO):} \quad $\square$ \textbf{NO AUTORIZADA} \quad $\square$ \textbf{SÍ (Tras pre-power físico)} \hfill \textbf{Firma Docente:} \rule{3cm}{0.4pt}
\end{tcolorbox}

\end{document}
